\documentclass[12pt]{article}
\usepackage{amsmath,amssymb,amsfonts,amsthm}
\usepackage[margin=1in]{geometry}

\begin{document}

\begin{center}
{\Large\bfseries Chapter 9, Problem 19}\\[6pt]
Byron \& Fuller, \textit{Mathematics of Classical and Quantum Physics}
\end{center}

\bigskip

\textbf{Problem.} Assuming that $\lambda$ is sufficiently small, expand $[D(\lambda)]^{-1}$ in powers of $\lambda$ and show that through order $\lambda^2$, $R(x,y;\lambda)$ is equal to the Neumann series expansion of the resolvent through order $\lambda^2$.

\bigskip
\textbf{Solution.}

The Fredholm resolvent is $R(x,y;\lambda) = N(x,y;\lambda)/D(\lambda)$.

\textbf{Neumann series:} The resolvent can also be obtained by iterating $f = g + \lambda Kf$:
\[
R(x,y;\lambda) = k(x,y) + \lambda\int k(x,t)k(t,y)\,dt + \lambda^2\int\!\!\int k(x,t_1)k(t_1,t_2)k(t_2,y)\,dt_1\,dt_2 + \cdots
\]
i.e., $R = K + \lambda K^2 + \lambda^2 K^3 + \cdots$ (in operator notation).

\textbf{Fredholm series:} From Problem~9.13:
\[
D(\lambda) = 1 - \lambda c_1 + \frac{\lambda^2}{2}c_2 - \cdots, \quad \text{where } c_1 = \int k(t,t)\,dt,
\]
\[
c_2 = \int\!\!\int\left[k(x_1,x_1)k(x_2,x_2) - k(x_1,x_2)k(x_2,x_1)\right]dx_1\,dx_2 = c_1^2 - \int\!\!\int |k(x_1,x_2)|^2\,dx_1\,dx_2.
\]

And:
\[
N(x,y;\lambda) = k(x,y) - \lambda\int\left[k(x,y)k(t,t) - k(x,t)k(t,y)\right]dt + O(\lambda^2).
\]
\[
= k(x,y) - \lambda\left[c_1 k(x,y) - k_2(x,y)\right] + O(\lambda^2),
\]
where $k_2(x,y) = \int k(x,t)k(t,y)\,dt$ is the iterated kernel.

\textbf{Expanding $D^{-1}$:}
\[
\frac{1}{D(\lambda)} = \frac{1}{1 - \lambda c_1 + O(\lambda^2)} = 1 + \lambda c_1 + \lambda^2 c_1^2 + O(\lambda^3) - \frac{\lambda^2}{2}c_2 + O(\lambda^3).
\]
More precisely:
\[
\frac{1}{D(\lambda)} = 1 + \lambda c_1 + \lambda^2\!\left(c_1^2 - \frac{c_2}{2}\right) + O(\lambda^3).
\]

\textbf{Computing $R = N/D$ through $O(\lambda^2)$:}
\begin{align}
R &= \left[k - \lambda(c_1 k - k_2) + O(\lambda^2)\right]\cdot\left[1 + \lambda c_1 + O(\lambda^2)\right] \\
&= k + \lambda c_1 k - \lambda c_1 k + \lambda k_2 + O(\lambda^2) \\
&= k + \lambda k_2 + O(\lambda^2).
\end{align}

At order $\lambda^2$: collecting all terms carefully:
\[
R = k + \lambda k_2 + \lambda^2 k_3 + O(\lambda^3),
\]
where $k_3(x,y) = \int\!\!\int k(x,t_1)k(t_1,t_2)k(t_2,y)\,dt_1\,dt_2$ is the third iterated kernel.

This is exactly the Neumann series through order $\lambda^2$:
\[
\boxed{R(x,y;\lambda) = k(x,y) + \lambda k_2(x,y) + \lambda^2 k_3(x,y) + O(\lambda^3)}.
\]

The agreement holds to all orders when $|\lambda| < 1/\|K\|$ (the radius of convergence of the Neumann series), confirming that the Fredholm and Neumann approaches give the same resolvent. \hfill $\blacksquare$

\end{document}
